% =====================================================================
% REUSABLE PROMPT (for future proposals in this style)
% ---------------------------------------------------------------------
% Copy everything below the next line into ChatGPT / Cursor / etc. when
% you want to regenerate a similar proposal for another iOS project.
%
% PROMPT START >>>>
% You are an academic writing assistant helping an NIBM BSc (Hons)
% Computing student write a coursework proposal for an iOS Application
% Development module. Write the proposal in LaTeX, around 1000 words,
% with a natural, confident, first-person voice — like a real student
% explaining their work, not a textbook. Avoid filler phrases such as
% "shall be able to". Keep sentences short. Use these sections:
%   1. Introduction
%   2. Problem Statement
%   3. Solution Overview (split into Mom side / Baby side / Midwife
%      side / Admin side, where applicable)
%   4. Mandatory iOS Features (Push notifications, Core Data,
%      Face ID / Touch ID, etc.)
%   5. Advanced iOS Feature (e.g., MapKit, WidgetKit)
%   6. Technologies and Architecture (MVVM, SwiftUI, Supabase, etc.)
%   7. Implementation Timeline (4-week table)
%   8. Testing Strategy
%   9. Risks and Ethics
%   10. Conclusion
%   11. References (APA style, Apple + framework docs)
% Use a clean title page with student name / ID / module placeholders.
% Inputs you should adapt the proposal to:
%   - App name and short tagline
%   - User roles (mom, baby, midwife, admin, etc.)
%   - Core features list
%   - Tech stack (SwiftUI, Supabase, Core Data, MapKit, WidgetKit,
%     LocalAuthentication, UserNotifications, etc.)
% Output: a single self-contained .tex file ready to compile.
% <<< PROMPT END
% =====================================================================

\documentclass[12pt,a4paper]{article}

\usepackage[margin=1in]{geometry}
\usepackage{setspace}
\usepackage{array}
\usepackage{longtable}
\usepackage{hyperref}
\usepackage{titlesec}

\onehalfspacing
\titleformat{\section}{\large\bfseries}{\thesection.}{0.5em}{}
\titleformat{\subsection}{\normalsize\bfseries}{\thesubsection.}{0.5em}{}

\begin{document}

\begin{titlepage}
    \centering
    \vspace*{1cm}
    {\Large \textbf{NIBM BSc (Hons) Computing}}\\
    \vspace{0.2cm}
    {\large iOS Application Development -- Coursework Proposal}\\
    \vspace{2cm}
    {\Huge \textbf{Prega Boo}}\\
    \vspace{0.4cm}
    {\large \textit{(originally proposed as ``MomCare'')}}\\
    \vspace{0.6cm}
    {\Large A Mother, Baby and Midwife Health Companion for iOS}\\
    \vspace{2cm}
    \begin{flushleft}
        \textbf{Student Name:} \underline{\hspace{8cm}}\\[0.4cm]
        \textbf{Student ID:} \underline{\hspace{8cm}}\\[0.4cm]
        \textbf{Module:} iOS Application Development\\[0.4cm]
        \textbf{Submission Type:} Project Proposal\\[0.4cm]
        \textbf{Tech Stack:} Swift, SwiftUI, Supabase, Core Data, MapKit, WidgetKit\\
    \end{flushleft}
    \vfill
    {\large \today}
\end{titlepage}

\section{Introduction}
Prega Boo is the iOS application I built for my coursework, and it grew out of a simple observation: in Sri Lanka, mothers and midwives still rely heavily on paper pregnancy books to track health information. Those books get torn, lost, or left at home on the day of a clinic visit. There is no single mobile app in the local context that connects mothers, babies and midwives in one place. Prega Boo (originally proposed as ``MomCare'') is my attempt to fix that gap with a clean, modern iOS app that treats pregnancy and early childcare as one continuous journey rather than two separate problems.

The app is written in Swift and SwiftUI, follows MVVM, and uses Supabase as the cloud backend. It supports four user roles -- expecting mothers, babies (managed under their mother's account), midwives and an admin account -- so the same install can be configured for different real-world users.

\section{Problem Statement}
Maternal health information in Sri Lanka is mostly captured on paper. That makes three things hard. First, mothers forget appointments and vaccinations because reminders depend on memory. Second, midwives cannot quickly check a mother's history before a home visit. Third, there is no reliable way to share notes between a mother and the midwife who looks after her. The result is missed clinic visits, incomplete vaccine schedules and avoidable confusion. A focused mobile app, built around real maternal-care workflows, can remove most of these issues.

\section{Solution Overview}
Prega Boo splits the experience into four sides so each user only sees what is relevant.

\textbf{Mom side.} After signing in (Google or email), a mother sees a calm dashboard that surfaces pregnancy week, daily care tips and quick links to track weight, mood, kicks, contractions and water intake. She can keep a private journal, view her own and her baby's profile, see vaccine and growth records added by her midwife, manage clinic appointments and find nearby hospitals through MapKit. A WidgetKit home-screen widget shows her name, district and baby count, and tapping it deep-links straight into the Mom \& Baby Details screen.

\textbf{Baby side.} Each mother can register one or more babies. The baby profile stores name, birth date, weight, photo and growth history. Vaccine schedules are built in: the midwife marks doses as given, and the mother sees the next due date.

\textbf{Midwife side.} Midwives log in to a list of mothers in their district. Tapping a mother opens her profile, baby cards, growth records, vaccine progress and a private notes section the midwife can write to. Notes appear instantly on the mother's side as a journal entry, so advice is never lost in a paper margin again.

\textbf{Admin side.} The admin can approve and manage midwife accounts, including profile photos and district assignments, all gated by Supabase Row Level Security policies.

\section{Mandatory iOS Features}
\textbf{Push and local notifications.} The app schedules local reminders for clinic appointments, vaccines and daily care nudges using the UserNotifications framework, and asks for permission politely on first launch.\\
\textbf{Core Data.} Selected data such as cached profiles and recent tracker entries are stored in Core Data so the app keeps working when the network is patchy. Cloud writes are layered on top through Supabase repositories.\\
\textbf{Face ID / Touch ID.} The app supports an optional PIN and biometric unlock through the LocalAuthentication framework. Face ID is auto-prompted when the lock screen appears (the way Apple Notes and banking apps do it), with PIN as a backup. Either method can be enabled on its own.

\section{Advanced iOS Features}
The two advanced features I leaned on are \textbf{MapKit} and \textbf{WidgetKit}. MapKit powers the Care Finder screen, where mothers see hospitals, clinics and maternal centres around them, with directions handed off to Apple Maps. WidgetKit powers a small and medium home-screen widget; tapping it triggers a custom URL scheme (\texttt{cw.prega-boo://widget-nav/mom-baby-details}) that a deep-link router translates into in-app navigation, even when the app is opened from a cold launch behind the lock screen.

\section{Technologies and Architecture}
The app follows the MVVM pattern. SwiftUI views own UI state, controllers prepare display models, and a repository layer wraps Supabase, Core Data, Keychain and Apple frameworks. This separation made it easy to swap implementations during development, add unit tests for tracker and PIN logic, and keep large screens readable.

The full stack is: Swift, SwiftUI, Xcode, Supabase (auth + Postgres + Storage + Row Level Security), Core Data, MapKit, WidgetKit, UserNotifications, LocalAuthentication, Keychain Services and GitHub for version control with branch-per-feature commits.

\section{Implementation Timeline}
\begin{longtable}{|p{2.2cm}|p{10.8cm}|}
\hline
\textbf{Week} & \textbf{Planned Work} \\
\hline
Week 1 & Finalise concept, wireframes and Supabase schema. Set up Xcode project, MVVM folders, GitHub. Build onboarding, welcome and login flows. \\
\hline
Week 2 & Mom dashboard, profile, pregnancy tracker, journal notes, water/weight/kicks/contractions, Core Data, push notifications. \\
\hline
Week 3 & Midwife and admin sides: mom list, baby cards, growth and vaccine records, midwife notes synced to mom, admin midwife management. \\
\hline
Week 4 & MapKit Care Finder, WidgetKit, Face ID/PIN polish, accessibility pass, manual + unit testing, bug fixes, final report. \\
\hline
\end{longtable}

\section{Testing Strategy}
I unit-tested the parts that hurt most when wrong: PIN hashing, kick/contraction calculations, date formatting and reminder scheduling. UI flows -- registration, sign-in, appointment creation, midwife note round-trip and widget deep-link -- were verified manually on iPhone simulators in light and dark mode, with Dynamic Type set to its larger sizes. Supabase Row Level Security policies were tested by signing in as different roles and confirming each could only read or write what they were meant to.

\section{Risks and Ethics}
Health data is sensitive, so the app collects only what is needed and stores PIN values as Keychain hashes, never as plain text. Supabase tokens live in Keychain too. Network errors are caught and shown as friendly messages instead of swallowed. The app does not give medical diagnosis -- tips are presented as supportive content with a clear suggestion to consult a healthcare professional.

\section{Conclusion}
Prega Boo started as a simple ``MomCare'' idea and ended up as a four-role iOS application that ties together SwiftUI design, MVVM architecture, Supabase, Core Data, MapKit, WidgetKit, push notifications and Face ID. It solves a real local problem -- replacing a paper pregnancy book with a digital companion that mothers, babies and midwives can all rely on -- while keeping the codebase modular enough to grow into a full maternal-care platform.

\section{References}
\begin{itemize}
    \item Apple. (2026). \textit{Human Interface Guidelines}. Apple Developer. \url{https://developer.apple.com/design/human-interface-guidelines/}
    \item Apple. (2026). \textit{SwiftUI documentation}. \url{https://developer.apple.com/documentation/swiftui}
    \item Apple. (2026). \textit{WidgetKit}. \url{https://developer.apple.com/documentation/widgetkit}
    \item Apple. (2026). \textit{LocalAuthentication}. \url{https://developer.apple.com/documentation/localauthentication}
    \item Apple. (2026). \textit{UserNotifications}. \url{https://developer.apple.com/documentation/usernotifications}
    \item Apple. (2026). \textit{MapKit}. \url{https://developer.apple.com/documentation/mapkit}
    \item Supabase. (2026). \textit{Supabase documentation}. \url{https://supabase.com/docs}
\end{itemize}

\end{document}
